\documentclass[10pt,letterpaper]{article}
\usepackage[utf8]{utf8}
\usepackage[margin=0.5in]{geometry}
\usepackage{hyperref}
\usepackage{enumitem}

\pagestyle{empty}
\hypersetup{colorlinks=true, linkcolor=black, urlcolor=blue}

\begin{document}
\fontfamily{phv}\selectfont

\begin{center}
    {\huge \textbf{ASHISH SURESH}}\\
    \vspace{1mm}
    \textbf{Embedded C++ Engineer}\\
    Goregaon West, Mumbai, Maharashtra | +91 8097177182 | \href{mailto:ashusresh98@gmail.com}{ashusresh98@gmail.com}\\
    \href{https://github.com/ashu304-ops/project1}{GitHub: github.com/ashu304-ops/project1}
\end{center}

\vspace{-2mm}

\section*{Profile}
Mechanical Engineer with 3 years of professional engineering experience transitioning into Embedded C++ development. Strong C++ foundation with hands-on project work in embedded C++, STM32 concepts, FreeRTOS, and Raspberry Pi software simulation. Seeking an entry-level/junior Embedded C++ or Firmware Engineer opportunity.

\vspace{-2mm}

\section*{Technical Skills}
\begin{itemize}[leftmargin=*, itemsep=2pt, parsep=0pt]
    \item \textbf{C++:} Modern C++, OOP, STL, Templates, RAII, Smart Pointers, C++17, FreeRTOS, STM32, Pointers, UART, SPI, I2C, Memory Management
    \item \textbf{Embedded \& Protocols:} C, STM32 Architecture, GPIO, ADC, PWM, Timers, Interrupts, UART, SPI, I2C
    \item \textbf{RTOS \& OS:} FreeRTOS (Tasks, Scheduling, Queues, Mutex, Semaphores, Software Timers), Linux Fundamentals, IPC
    \item \textbf{Development Tools:} CMake, GCC/GDB, Git, GitHub, Raspberry Pi Software Simulation
    \item \textbf{Concepts \& Architecture:} SOLID Principles, Design Patterns, Debugging, Testing, Embedded Linux Integration
\end{itemize}

\vspace{-2mm}

\section*{Embedded Engineering Projects}
\textbf{Embedded Vehicle Suspension \& Safety Controller} \hfill \href{https://github.com/ashu304-ops/project1}{GitHub Repository}\\
\textit{C++ | STM32 | FreeRTOS | Raspberry Pi Simulation}
\begin{itemize}[leftmargin=*, itemsep=2pt, parsep=0pt]
    \item Developed an Embedded C++ vehicle control system with simulated hardware sensors, actuators, and safety mechanisms.
    \item Applied Object-Oriented Design, RAII, and STL to design modular C++ firmware components for real-time operations.
    \item Implemented FreeRTOS tasks, message queues, and synchronization primitives for sensor acquisition and safety controls.
    \item Configured virtual peripherals for STM32 concepts including GPIO, ADC, PWM, timers, interrupts, UART, SPI, and I2C.
    \item Built a Raspberry Pi software simulation environment to emulate hardware interactions, sensor failures, and thermal faults.
    \item Structured system architecture to prepare for future Embedded Linux integration and real-time visual monitoring.
\end{itemize}

\vspace{-2mm}

\section*{Professional Experience}
\textbf{Dipesh Engineering Pvt. Ltd.} -- Design Engineer \hfill \textbf{Mumbai, India}\\
\textit{Professional Engineering (3 Years)}
\begin{itemize}[leftmargin=*, itemsep=2pt, parsep=0pt]
    \item Designed complex engineering systems according to thermal, mechanical, and manufacturing specifications.
    \item Executed structured engineering analysis, parametric calculations, and technical optimization for clients.
    \item Collaborated with cross-functional manufacturing and engineering teams to resolve technical design edge cases.
\end{itemize}

\vspace{-2mm}

\section*{Education}
\textbf{Bachelor of Engineering} -- Mechanical Engineering \hfill \textbf{Mumbai, India}\\
Thakur College of Engineering and Technology

\end{document}
